\documentclass[11pt]{article}
\usepackage{amsmath,amssymb}
\usepackage[hidelinks]{hyperref}
\title{A truthful heading is not a commitment to a future bay}
\author{}\date{}
\begin{document}\maketitle

\paragraph{Scope and interface.}
Continue the scouting experiment with an old carrier permitted to stay
at its bay or start a cross-bay maneuver. The new carrier must reach the
open bay by tick 8 without blocking the old carrier. A separate required
job still runs. The scout's costly step replaces an optional pickup,
and is needed to see the affected cells. The local input comes from the
pinned RWARE dictionary observation, not a planner-supplied fault label.
It reports presence and orientation in the local crop. The legal schedule
and actor commitments are authored; pinned selected RWARE commands and
SubT delivery and charging supply the moving and transport parts.

\paragraph{Three legal continuations and two collapses.}
In the stationary branch the observed old carrier remains pointed
upward at its original bay. In the crossing branch it actually turns
toward the other bay before the observation, and later executes two
authorized forward steps. The turn has changed heading but not
position at the comparison cut. Thus a completely accurate placement
read agrees between the two cases, while the open bay at arrival
differs. The new carrier must choose before the crossing steps.

A scope-qualified heading read distinguishes the two if the
turned carrier is known to complete the authorized crossing.
With a permissible withdrawal \emph{after the read}, both the local crop
and any delivered report derived solely from it agree with the crossing
case. The two actual final bay occupancies then disagree. There is
no branch policy on that reading alone that guarantees both. This is a
claim about making the branch before new evidence, not about every
future data acquisition. A fresh proof of the applicable commitment,
a lawful time-bounded restriction on withdrawal, or a later useful
revision with enough route time would change the premises.

\paragraph{Execution and costs.}
With a free common prefix, the carrier checks the cell and moves
before the tick-4 cue. With a restricted prefix it uses a tick-3 cue
before its first turn. The two supplied timing cuts follow the
previous scout-frontier result. The witness preserves the old
crossing instead of freezing it to make the forecast correct.
The probe loses a scheduled pickup, and a sent cue costs three
native TX bytes even if erased. Independent motion and the
independent radio call both advance. The comparison doesn't treat an
authorized withdrawal as completion of the earlier physical crossing;
it verifies the corresponding authorized terminal branch instead.

\paragraph{Evidence and residual.}
Ninety-six selected native comparisons span both initial bays,
three old programs, link erasure, two prefix permissions and four
policies. The heading rule meets the new deadline in all eight
delivered, stable-law cases. Extending it without qualification to
withdrawal fails in each delivered withdrawal case. Position-only
and unqualified heading each hit eight of 24 overall, but at different
histories: the former chooses incorrectly for a crossing; the latter
chooses incorrectly for a withdrawal. This balanced test count is not
a failure rate. A guard uncertain about withdrawal can still use
the non-turning case; on ambiguous turns it forgoes timely new service.
The exact direct cases agree, and Lean verifies the scoped prediction,
the equal-read/opposite-outcome pair, and the one-shot impossibility.
No endpoint-independent communication claim, demonstrated mechanism
for establishing scope, new optimization principle or held-out
mission is inferred.
\end{document}
